\documentclass{article}

\usepackage{xcolor}
\usepackage[utf8]{inputenc}
\usepackage{amsmath}
\usepackage{enumerate}
\usepackage{microtype}
\usepackage[top=1.5cm]{geometry} % controls page margins

\title{Examen Parcial 2 - Bonus}
\author{Estructuras de Datos \\ Facultad de Ingenierías - Universidad Tecnológica de Pereira \\ Ingeniería de Sistemas y Computación}
\date{Abril 29 2026}

\begin{document}
\maketitle
\thispagestyle{empty}

Considere el archivo \texttt{bonus.c} adjunto.

\vspace{0.3cm}

Debe corregir el archivo de manera que se comporte de acuerdo con lo esperado. Es decir, debe tomar un conjunto finito de valores enteros, almacenarlos en un montículo tipo \textit{minHeap}, y posteriormente insertar dichos valores, en ese orden, dentro de un árbol AVL.

\vspace{0.3cm}

Considere que, para este código, el montículo es creado con un tamaño fijo y se está trabajando con una estructura que no necesariamente debe permitir redimensionamiento dinámico (array).

\end{document}